\section{实验和 Linux 源码阅读路线}

算法书需要实验，是因为很多概念只有跑起来才会稳定。第一组实验是复杂度实验：分别写平方枚举、前缀和、哈希前缀计数，用不同 n 观察增长曲线。第二组实验是容器实验：比较 vector 顺序扫描、list 顺序扫描、unordered\_map 查询和 map 查询，在数据量变大后观察缓存局部性和常数差异。第三组实验是状态实验：把 BFS 的访问标记从入队改到出队，观察重复入队数量如何膨胀。

Linux 源码阅读可以围绕数据结构而不是模块盲逛。先看链表宏如何通过成员地址反推出外层结构体，再看红黑树如何把比较逻辑留给调用者，再看 XArray 或基数树如何服务于页缓存和 ID 分配。阅读目标不是背字段名，因为字段会随版本变化；目标是理解为什么内核倾向侵入式结构、为什么要避免频繁小对象分配、为什么并发读写需要锁或 RCU 这类同步机制。

把刷题题型和内核结构对应起来会很有帮助。LRU 缓存对应哈希表加双向链表；定时器或调度候选对应按时间或虚拟运行时间排序的有序结构；文件描述符和页缓存索引对应整数到对象的映射；网络包队列对应队列和优先队列；内存伙伴系统对应分层空闲链表。对应关系不是说内核用的就是力扣模板，而是说明数据结构选择永远由访问模式决定。

源码实验要小步进行。可以先写一个最小侵入式链表：业务对象中放 prev 和 next，通过容器\_of 思想从链表节点拿回业务对象。再写一个普通 std::list 版本，对比所有权和迭代器稳定性。然后写一个 LRU，分别用 std::list 加 unordered\_map 和手写双向链表实现。这个实验能把链表、哈希、对象生命周期和异常安全放在同一张图里。

性能实验要避免伪结论。先固定编译优化级别，预热数据，避免把输出打印放进计时代码；再多次运行取中位数；最后解释结果。vector 往往因为连续内存比 list 更快，即使两者同为线性扫描。unordered\_map 平均常数很快，但如果 key 构造很重，整体成本仍可能高。算法复杂度给上界，工程测量暴露常数、缓存和分配成本。

每章实验都应产出一页报告：问题是什么，暴力基线是什么，优化版本改变了哪一步，输入规模如何设置，结果是否符合复杂度预期，出现偏差的原因是什么。报告不追求华丽，追求可复现。以后复习时，能用自己的实验解释一个算法，比只背题解更可靠。

\begin{principlebox}{专题复盘要求}
读完本节后，不要只记结论。请把每个结论还原成一个可验证的问题：它解决了哪一步重复工作，依赖哪个题目条件，使用哪个容器或状态，边界条件如何破坏错误写法。
\end{principlebox}
